\documentclass[12pt]{article}
% AAU theme adapted from the preamble supplied by the project team.
\usepackage[a4paper,margin=1in]{geometry}
\usepackage{lmodern}
\usepackage[T1]{fontenc}
\usepackage{microtype}
\usepackage{graphicx}
\usepackage{float}
\usepackage{amsmath}
\usepackage[table]{xcolor}
\usepackage{titlesec}
\usepackage{fancyhdr}
\usepackage[most]{tcolorbox}
\PassOptionsToPackage{obeyspaces}{url}
\usepackage[hidelinks]{hyperref}
\usepackage{array}
\usepackage{tabularx}
\usepackage{enumitem}
\usepackage{booktabs}
% Vector UML artwork and multipage test tables.
\usepackage{longtable}
\usepackage{pdflscape}
\definecolor{aauRed}{RGB}{164,22,35}
\definecolor{aauGreen}{RGB}{0,122,88}
\definecolor{darkGray}{RGB}{60,60,60}
\definecolor{softGray}{RGB}{245,245,245}
\pagestyle{fancy}\fancyhf{}\fancyfoot[C]{\thepage}
\renewcommand{\headrulewidth}{0pt}\renewcommand{\footrulewidth}{0pt}
\titleformat{\section}{\normalfont\Large\bfseries\color{aauRed}}{\thesection}{0.6em}{}[{\color{aauRed!65}\titlerule[1.1pt]}]
\titlespacing*{\section}{0pt}{12pt}{8pt}
\titleformat{\subsection}{\normalfont\normalsize\bfseries\color{aauRed!85!black}}{\thesubsection}{0.6em}{}
\titlespacing*{\subsection}{0pt}{10pt}{4pt}
\newtcolorbox{objectivebox}{enhanced,colback=aauRed!5,colframe=aauRed!80!black,boxrule=0.3mm,arc=1mm,left=10pt,right=10pt,top=6pt,bottom=6pt,borderline west={2.5pt}{0pt}{aauRed}}
% Each short use/test case stays on one page, including its heading.
\newtcolorbox{contentbox}[1][]{enhanced,colback=softGray,colframe=black!35,boxrule=0.25mm,arc=2mm,left=9pt,right=9pt,top=7pt,bottom=7pt,before skip=9pt,after skip=9pt,title={\textbf{#1}},fonttitle=\bfseries,coltitle=darkGray,colbacktitle=black!8}
\newcommand{\blankline}{\rule{\linewidth}{0.35pt}}
\newcommand{\insertfigure}[3]{\begin{figure}[H]\centering\fbox{\parbox[c][#2][c]{0.88\textwidth}{\centering\textcolor{black!50}{\Large #1}}}\caption{#3}\end{figure}}
\newcommand{\screen}[3][0.98]{\begin{figure}[H]\centering\includegraphics[width=#1\linewidth]{../assets/#2}\caption{#3}\end{figure}}
% Give wide UML models a dedicated landscape page with readable labels.
\newenvironment{umlpage}{%
\newgeometry{margin=0.65in}\begin{landscape}\pagestyle{empty}%
\AddToHook{shipout/foreground}[arena/umlfooter]{%
\begin{tikzpicture}[remember picture,overlay]
\node[rotate=90,anchor=south,inner sep=0pt] at ([xshift=-23pt]current page.east)
{\normalfont\normalsize\thepage};
\end{tikzpicture}}%
}{\end{landscape}\RemoveFromHook{shipout/foreground}[arena/umlfooter]\pagestyle{fancy}\restoregeometry}
\newcommand{\umlfigure}[2]{\begin{figure}[H]\centering
\includegraphics[width=\linewidth,height=13cm,keepaspectratio]{../assets/uml/#1.pdf}
\caption{#2}\end{figure}}
\providecommand{\reportdate}{10 November 2026}
% Two intentionally sized figures fit together with a heading on an A4 page.
\newcommand{\screenpair}[5][0.30]{%
\begin{figure}[H]\centering
\includegraphics[width=0.96\linewidth,height=#1\textheight,keepaspectratio]{../assets/#2}\caption{#3}
\vspace{8pt}
\includegraphics[width=0.96\linewidth,height=#1\textheight,keepaspectratio]{../assets/#4}\caption{#5}
\end{figure}}
\setlength{\parindent}{0pt}\setlength{\parskip}{5pt}
\setlist[itemize]{leftmargin=1.5em,itemsep=2pt,topsep=3pt}
\setlist[enumerate]{leftmargin=1.7em,itemsep=3pt,topsep=4pt}
\setlength{\emergencystretch}{2em}
\widowpenalty=10000\clubpenalty=10000\displaywidowpenalty=10000
\brokenpenalty=10000
\raggedbottom
\renewcommand{\arraystretch}{1.16}
\setlength{\abovecaptionskip}{7pt}
\setlength{\belowcaptionskip}{4pt}
\setlength{\intextsep}{12pt plus 2pt minus 2pt}
\newcolumntype{Y}{>{\raggedright\arraybackslash}X}
\BeforeBeginEnvironment{tabularx}{\par\noindent}
\AfterEndEnvironment{tabularx}{\par}
\newcommand{\code}[1]{{\small\expandafter\nolinkurl\expandafter{\detokenize{#1}}}}
\newcommand{\command}[1]{\mbox{\small\ttfamily\detokenize{#1}}}
\newcommand{\projecttitle}[2]{%
\hypersetup{pageanchor=false}\begin{titlepage}\centering\vspace*{-0.45cm}
\noindent\begin{minipage}[c]{0.66\textwidth}\raggedright
{\Large\bfseries Al Ain University (AAU)\par}
\end{minipage}\hfill
\begin{minipage}[c]{0.28\textwidth}\raggedleft
% The supplied PNG is unchanged; trim removes only its transparent margins.
\includegraphics[width=3.1cm,trim=1731bp 871bp 1735bp 878bp,clip]{../../aau.png}
\end{minipage}\par
\vspace{0.45cm}
{\scshape\large Object Oriented Analysis and Design\par}
\vspace{0.10cm}
{\normalsize System Analysis \& Design\par}
\vspace{0.10cm}
{\small Course codes: 0103430 \;|\; 0102331\par}
\vspace{0.3cm}\noindent\rule{\linewidth}{0.5mm}\par\vspace{0.3cm}
{\huge\bfseries Arena\par}
\vspace{0.18cm}{\Large\color{aauRed}#1\par}
\vspace{0.15cm}{\normalsize\color{darkGray}#2\par}
\vspace{0.12cm}{\color{aauRed}\rule{\linewidth}{1.2pt}\par}
\vspace{0.35cm}{\normalsize\bfseries Prepared by\par}\vspace{0.15cm}
\begin{tabular}{@{}lr@{}}\textbf{Name}&\textbf{Student ID}\\\midrule
Khaled Alkhawaldeh&202410091\\Abdulrahman Alnaqbi&202340042\\Yousif Alhammadi&202310707\\Ahmed Alsamma&202320192\end{tabular}\par
\vspace{0.22cm}{\color{aauRed!55}\rule{0.6\textwidth}{1.2pt}\par}
\vspace{0.15cm}{\normalsize\textbf{Section:} 61D\par}
\vspace{0.15cm}{\normalsize\textbf{Submitted to:} Dr. Yazeed Ghadi\par}
\vspace{0.3cm}
\begin{tcolorbox}[colback=gray!5,colframe=black!70,title=\textbf{Project Summary},fonttitle=\bfseries,width=0.95\textwidth,boxrule=0.4mm,arc=3mm,left=7pt,right=7pt,top=4pt,bottom=4pt]
A Java desktop workspace for organizing campus competitions, coordinating teams, collecting submissions, evaluating work, and publishing results. Project, programming, and quiz scoring are separated from score-ranking, knockout, and league progression through object-oriented strategies. The application is developed for Al Ain University in the United Arab Emirates.
\end{tcolorbox}
\vfill{\normalsize Fall, 2026--2027\par}{\small Version 2.0 \;|\; \reportdate\par}
\end{titlepage}\pagenumbering{roman}\hypersetup{pageanchor=true}\tableofcontents\clearpage\pagenumbering{arabic}}

\begin{document}
\projecttitle{Testing Report}{Acceptance cases, execution evidence, and quality assessment}
\section{Introduction}
\subsection{Purpose of Testing}
Testing verifies that Arena implements the specified competition policies and preserves them when actions are invalid or storage fails. The main objective is to validate observable behavior and invariants rather than mirror individual lines of implementation.
\subsection{Testing Scope}
The suite covers authentication, account creation, ownership and role checks, competition phases, team membership, registrations, submissions, evaluations, rankings, export, knockout progression, password storage, file locking, restart persistence, corrupt input, and failed-write rollback. Network integrations, automatic grading, multiple clients, accessibility certification, and performance benchmarks are outside the executed scope.

\section{Test Strategy}
\subsection{Unit Testing}
Direct policy checks validate the 40/35/25 project formula, score bounds, password verification and salt uniqueness, and ranking outcomes. These are executed within the same dependency-free runner as the integration scenarios.
\subsection{Integration Testing}
ArenaService is exercised against a real Repository in a uniquely named build subdirectory. Most tests use an injected fixed UTC clock representing 29 October 2026; a separate clock places submission exactly at the deadline. The test dataset is isolated from the user's application database. Tests log in as distinct actors instead of passing forged role flags.

The persistence tests reopen the saved file, attempt a concurrent repository lock, supply an invalid serialized file, and deliberately obstruct the temporary save path. The failed-write case verifies that the service restores its prior team count rather than leaving an unsaved in-memory change.
\subsection{System Testing}
Complete application-service workflows cover project judging, programming results, and a knockout champion. A separate Swing suite performs 22 checks through actual controls on the event-dispatch thread, using an isolated repository. It enters credentials, clicks buttons, changes search/filter controls, creates a team, records a match winner, checks the live project total and saves its review, cancels a dialog, and checks layout at the minimum window size. These are programmatic control interactions, not operating-system mouse automation. Thirteen actual Swing screens and dialogs were also rendered for visual inspection. A human usability study was not performed. The test runner returns a nonzero process exit code if any check fails.

\clearpage\section{Test Cases}
\subsection{Authentication Test Cases}
\begin{contentbox}[TC-02: Incorrect credentials clear the session]
\textbf{Requirement:} FR-01. \textbf{Precondition:} a valid organizer session exists.\\
\textbf{Steps:} call login with the organizer email and an incorrect password; attempt to read the current user.\\
\textbf{Expected:} login reports incorrect credentials; subsequent session access requires sign-in.\\
\textbf{Actual:} both rejections occurred. \textbf{Status: Pass.}
\end{contentbox}
\begin{contentbox}[TC-42: Password change invalidates the old password]
\textbf{Requirement:} NFR-02 / FR-02. \textbf{Precondition:} student10 is authenticated.\\
\textbf{Steps:} change password after supplying the current password; sign out; try old password; try new password.\\
\textbf{Expected:} old credential fails and new credential succeeds.\\
\textbf{Actual:} expected behavior observed. \textbf{Status: Pass.}
\end{contentbox}
\subsection{Competition Management Test Cases}
\begin{contentbox}[TC-10: Another organizer cannot advance the competition]
\textbf{Requirement:} FR-03. \textbf{Precondition:} organizer1 owns a new competition.\\
\textbf{Steps:} sign in as organizer2 and request the next phase.\\
\textbf{Expected:} ownership rejection; the competition remains in Draft.\\
\textbf{Actual:} the owner-only guard rejected the request. \textbf{Status: Pass.}
\end{contentbox}
\begin{contentbox}[TC-28: Incomplete reviews block publication]
\textbf{Requirement:} FR-08. \textbf{Precondition:} a two-team project is in Judging with only a subset of assigned reviews.\\
\textbf{Steps:} sign in as the owner and request publication.\\
\textbf{Expected:} publication is rejected until every assigned judge evaluates every entry.\\
\textbf{Actual:} incomplete-review guard rejected publication. \textbf{Status: Pass.}
\end{contentbox}
\subsection{Registration Test Cases}
\begin{contentbox}[TC-15: Duplicate competition registration]
\textbf{Requirement:} FR-04. \textbf{Precondition:} Team Alpha is registered and its captain is signed in.\\
\textbf{Steps:} request registration of Team Alpha in the same competition again.\\
\textbf{Expected:} duplicate entry is rejected.\\
\textbf{Actual:} the service reported that the team was already registered. \textbf{Status: Pass.}
\end{contentbox}
\begin{contentbox}[TC-18: A student cannot enter twice through different teams]
\textbf{Requirement:} FR-04. \textbf{Precondition:} student1 belongs to registered Team Alpha.\\
\textbf{Steps:} student1 creates another team and tries to register it in the same event.\\
\textbf{Expected:} overlapping membership blocks registration.\\
\textbf{Actual:} the service rejected the member overlap. \textbf{Status: Pass.}
\end{contentbox}
\subsection{Submission Test Cases}
\begin{contentbox}[TC-22: A valid resubmission replaces earlier work]
\textbf{Requirement:} FR-06. \textbf{Precondition:} active project, registered captain, future deadline.\\
\textbf{Steps:} submit a project; submit a revised title, description, and link.\\
\textbf{Expected:} one current submission exists, with the revised title.\\
\textbf{Actual:} submission count was one and the title matched the revision. \textbf{Status: Pass.}
\end{contentbox}
\begin{contentbox}[TC-47: Exact-deadline submission is too late]
\textbf{Requirement:} FR-06. \textbf{Precondition:} active competition with registered teams.\\
\textbf{Steps:} reopen the service with the clock fixed exactly at the deadline; attempt submission.\\
\textbf{Expected:} reject, because the deadline is an exclusive upper bound.\\
\textbf{Actual:} deadline rejection occurred. \textbf{Status: Pass.}
\end{contentbox}
\subsection{Evaluation and Leaderboard Test Cases}
\begin{contentbox}[TC-35: Weighted project score]
\textbf{Requirement:} FR-07 / project policy. \textbf{Precondition:} an evaluation contains innovation 90, execution 80, presentation 70.\\
\textbf{Steps:} pass the evaluation to ProjectRules.score.\\
\textbf{Expected:} $90(0.4)+80(0.35)+70(0.25)=81.5$.\\
\textbf{Actual:} computed score equaled 81.5 within numerical tolerance. \textbf{Status: Pass.}
\end{contentbox}
\begin{contentbox}[TC-38 and TC-41: Knockout progression to champion]
\textbf{Requirement:} FR-10. \textbf{Precondition:} four-team active bracket and assigned judge3.\\
\textbf{Steps:} decide both semifinals; inspect the generated final; decide the final; certify and publish as owner.\\
\textbf{Expected:} exactly one final is created; champion is first with two wins.\\
\textbf{Actual:} three total matches existed and Team Alpha led with two wins after publication. \textbf{Status: Pass.}
\end{contentbox}
\begin{contentbox}[TC-48: Persistence failure rolls back the update]
\textbf{Requirement:} NFR-04. \textbf{Precondition:} writable isolated test repository.\\
\textbf{Steps:} obstruct its temporary-file path with a nonempty directory; attempt team creation; compare team count; remove only the test obstruction.\\
\textbf{Expected:} save error and unchanged in-memory team count.\\
\textbf{Actual:} rollback preserved the team count. \textbf{Status: Pass.}
\end{contentbox}

\section{Test Results}
\subsection{Test Execution Summary}
Execution date: 10 November 2026. Platform: Windows. Build: portable JDK 21, Java 17 target. Result: \textbf{70 existing service checks, 28 Arena regression checks, and 22 Swing UI checks passed; 0 failed}. All application suites can be run with:
\begin{verbatim}
.\build-app.ps1 -Test -UiTest
\end{verbatim}
Maven integration uses three JUnit wrapper methods. Default \command{mvn test} runs the 70 original service checks and 28 Arena regression checks, and skips the Swing wrapper. \command{mvn test -Pui-tests} also runs all 22 Swing checks. Surefire reports three wrapper tests; these are not 120 independent JUnit methods. A graphical session is required for Swing. Both native runners and Maven fail if a check fails. CSV files and timestamped summaries are preserved in \code{reports/test-evidence}.

The optional browser runner executed \textbf{17 additional checks, all passing}, covering prototype branding, discovery, project submission and scoring, knockout progression, quiz ties, league results, creation-preview compatibility, and four viewport widths. It uses isolated in-memory browser data. Browser checks do not substitute for Java service or persistence tests.

\begin{verbatim}
mvn clean compile
mvn test -Pui-tests
node scripts/test-prototype.cjs --browser PATH_TO_BROWSER
\end{verbatim}
The browser runner requires Node.js, Playwright, and an installed supported browser; it is optional tooling, not an application runtime dependency.

\subsection{Passed and Failed Tests}
Each row below records an executed check. The CSV preserves the same IDs and outcomes; the Java source contains the exact setup, steps, values, and assertions. Scenarios share a deliberate lifecycle fixture, so run the whole suite rather than assuming each row is independently executable.
\small\begin{longtable}{p{0.12\linewidth}p{0.66\linewidth}p{0.1\linewidth}}\toprule ID & Requirement / executed check & Result\\\midrule\endhead
TC-01 & FR-01: valid authentication & PASS\\
TC-02 & FR-01: wrong password clears session & PASS\\
TC-03 & FR-01: unknown account rejected & PASS\\
TC-04 & FR-02: student cannot create accounts & PASS\\
TC-05 & FR-02: administrator creates an account & PASS\\
TC-06 & FR-02: duplicate email is case insensitive & PASS\\
TC-07 & FR-02: short passwords rejected & PASS\\
TC-08 & FR-03: student cannot create competitions & PASS\\
TC-09 & FR-03: new competition is a draft & PASS\\
TC-10 & FR-03: only owning organizer manages competition & PASS\\
TC-11 & FR-04: draft registration rejected & PASS\\
TC-12 & FR-03: rule edits blocked after opening & PASS\\
TC-13 & FR-04: at least two entries required & PASS\\
TC-14 & FR-04: captain registers team & PASS\\
TC-15 & FR-04: duplicate registration rejected & PASS\\
TC-16 & FR-04: noncaptain registration rejected & PASS\\
TC-17 & FR-05: registered roster is locked & PASS\\
TC-18 & FR-04: same student cannot enter twice & PASS\\
TC-19 & FR-04: withdrawal during registration & PASS\\
TC-20 & FR-06: missing submissions block judging & PASS\\
TC-21 & FR-06: non-HTTP submission URL rejected & PASS\\
TC-22 & FR-06: submission may be replaced before judging & PASS\\
TC-23 & FR-06: student sees only own team's submissions & PASS\\
TC-24 & FR-07: scores rejected outside judging & PASS\\
TC-25 & FR-07: unassigned judge rejected & PASS\\
TC-26 & FR-07: scores outside rubric rejected & PASS\\
TC-27 & FR-07: replacing judge score does not duplicate it & PASS\\
TC-28 & FR-08: unfinished reviews block publication & PASS\\
TC-29 & FR-08: student cannot read unpublished leaderboard & PASS\\
TC-30 & FR-08: equal project averages share rank & PASS\\
TC-31 & FR-07: published scores cannot be modified & PASS\\
TC-32 & FR-08: standings export has header and rows & PASS\\
TC-33 & FR-03: archive is terminal & PASS\\
TC-34 & FR-09: programming ranks solved then penalty & PASS\\
TC-35 & FR-07: weighted project formula & PASS\\
TC-36 & FR-10: only assigned judge records bracket results & PASS\\
TC-37 & FR-10: winner must play in the selected match & PASS\\
TC-38 & FR-10: semifinals create exactly one final & PASS\\
TC-39 & FR-10: finalized match cannot be overwritten & PASS\\
TC-40 & FR-10: unfinished bracket blocks certification & PASS\\
TC-41 & FR-10: champion is published after final & PASS\\
TC-42 & NFR-02: password changes invalidate old password & PASS\\
TC-43 & NFR-02: salted hashes differ for equal passwords & PASS\\
TC-44 & NFR-04: concurrent data access is locked & PASS\\
TC-45 & FR-11: successful changes are audited & PASS\\
TC-46 & FR-11: students cannot read administrative audit & PASS\\
TC-47 & FR-06: submission at deadline is rejected & PASS\\
TC-48 & NFR-04: failed disk write rolls back domain changes & PASS\\
TC-49 & FR-03: knockout capacity cannot exceed the supported bracket limit & PASS\\
TC-50 & FR-03: past deadlines rejected & PASS\\
TC-51 & FR-03: nonjudge cannot be assigned as judge & PASS\\
TC-52 & FR-09: one official programming result per team & PASS\\
TC-53 & FR-09: programming results complete publication workflow & PASS\\
TC-54 & FR-04: oversized teams rejected & PASS\\
TC-55 & FR-05: student can create an unlocked team & PASS\\
TC-56 & FR-05: unlocked membership changes persist & PASS\\
TC-57 & FR-08: CSV formula injection escaped & PASS\\
TC-58 & NFR-04: restart preserves accounts and published champion & PASS\\
TC-59 & NFR-04: corrupt data is not silently reset & PASS\\
TC-60 & SEC: account APIs return credential-free summaries & PASS\\
TC-61 & SEC: account directory is administrator-only & PASS\\
TC-62 & SEC: null credentials reject and clear the session & PASS\\
TC-63 & FR-02: missing account role leaves accounts unchanged & PASS\\
TC-64 & FR-06: relative URL rejects without a null-pointer error & PASS\\
TC-65 & SEC: malformed and excessive password derivation parameters reject & PASS\\
TC-66 & FR-08: multiline formula cells are escaped in exported CSV & PASS\\
TC-67 & NFR-04: oversized snapshot preserves the last readable file & PASS\\
TC-68 & NFR-04: closed repository rejects writes and permits repeated close & PASS\\
TC-69 & NFR-04: pre-refactor version-one database remains readable & PASS\\
TC-70 & FR-10: a null winner rejects without changing the bracket & PASS\\
\bottomrule\end{longtable}\normalsize

\subsection{Arena Type and Tournament Regression Results}
These checks include all seven valid combinations, rejection of unsupported configurations, quiz bounds and shared ranks, every bracket and league size from 2 to 16, byes, unique fixtures, all league tie-break criteria, authorization, lifecycle guards, result immutability, failed-save rollback, restart persistence, and both legacy fixtures. AR-11 and AR-15 loop over every supported team count; each loop remains one named check, not fifteen reported tests.
\small\begin{longtable}{p{0.12\linewidth}p{0.66\linewidth}p{0.1\linewidth}}\toprule ID & Arena regression check & Result\\\midrule\endhead
AR-01 & All seven supported type and tournament combinations save and reload & PASS\\
AR-02 & Unsupported combinations reject without adding records or audit entries & PASS\\
AR-03 & Null type or tournament and invalid quiz maxima reject & PASS\\
AR-04 & Students cannot create tournaments and other organizers cannot edit them & PASS\\
AR-05 & Draft type changes persist and registration freezes the configuration & PASS\\
AR-06 & Quiz ranking enters judging without submissions and rejects submitted work & PASS\\
AR-07 & Quiz points enforce lower upper and irrelevant-input boundaries & PASS\\
AR-08 & Quiz publication needs all results and equal top scores share first place & PASS\\
AR-09 & Quiz official result belongs to one assigned judge & PASS\\
AR-10 & Students cannot see provisional league tables or enter match scores & PASS\\
AR-11 & All 2 through 16 team knockout brackets finish with exactly n minus one played matches & PASS\\
AR-12 & Tournament entrants reject duplicate null too-small and oversized lists & PASS\\
AR-13 & Three-team quiz knockout records points advances byes and publishes a champion & PASS\\
AR-14 & Five-team general knockout preserves winner selection and rejects repeated results & PASS\\
AR-15 & All 2 through 16 team league schedules contain every pair once with no round collisions & PASS\\
AR-16 & League win draw loss counters and 3-1-0 points are accurate & PASS\\
AR-17 & League tied points use wins then score difference and scores for & PASS\\
AR-18 & Completely tied league results use case-insensitive names then stable IDs & PASS\\
AR-19 & League cannot certify early and completed fixtures publish and export full statistics & PASS\\
AR-20 & Quiz leagues accept draws within configured points and reject out-of-range scores & PASS\\
AR-21 & General league scoring rejects values above the supported limit & PASS\\
AR-22 & Failed match write rolls back the result new round and audit then succeeds on retry & PASS\\
AR-23 & Failed league write leaves standings and fixture unchanged & PASS\\
AR-24 & Quiz and league results survive a complete repository restart & PASS\\
AR-25 & Legacy bytes migrate without rewriting on read and save as schema two & PASS\\
AR-26 & Legacy active knockout continues to champion in Arena & PASS\\
AR-27 & Failed quiz write restores the prior evaluation & PASS\\
AR-28 & Match scores cannot change after publication and archived events cannot advance & PASS\\
\bottomrule\end{longtable}\normalsize

\subsection{Swing Interaction Results}
These 22 cases exercise real Swing controls and dialogs. The CSV results and timestamped summary in the evidence folder preserve their recorded outcomes.
\begingroup\renewcommand{\arraystretch}{1.05}\footnotesize\begin{longtable}{p{0.12\linewidth}p{0.66\linewidth}p{0.1\linewidth}}\toprule ID & Swing interaction check & Result\\\midrule\endhead
UI-01 & Sign-in controls and packaged university artwork are available & PASS\\
UI-02 & Organizer signs in through the real form and sees permitted navigation & PASS\\
UI-03 & Search handles no results and restores cards when cleared & PASS\\
UI-04 & Format filter and competition card navigate to the correct event & PASS\\
UI-05 & Bracket tab displays the seeded round & PASS\\
UI-06 & Results page shows published standings & PASS\\
UI-07 & Sign-out returns to an empty sign-in form & PASS\\
UI-08 & Student navigation hides organizer and judge functions & PASS\\
UI-09 & Create-team dialog saves through the service and refreshes the roster & PASS\\
UI-10 & Judging studio lists only the current judge's assignments & PASS\\
UI-11 & Match-winner dialog persists the selected team & PASS\\
UI-12 & Administrator opens the account directory & PASS\\
UI-13 & Canceling account creation leaves stored users unchanged & PASS\\
UI-14 & Competition cards remain within the viewport at minimum window size & PASS\\
UI-15 & Project scoring dialog calculates a live total and saves the review & PASS\\
UI-16 & User-supplied markup remains literal in names and table cells & PASS\\
UI-17 & Arena branding appears in window login and sidebar & PASS\\
UI-18 & Creation form filters compatible formats and reveals quiz settings & PASS\\
UI-19 & Creation form saves a quiz league with selected settings & PASS\\
UI-20 & Quiz scoring presents one points control and saves without submissions & PASS\\
UI-21 & League fixture dialog saves a draw and refreshes the results & PASS\\
UI-22 & Workspace remains within the viewport at 1100 by 720 & PASS\\
\bottomrule\end{longtable}\endgroup

\subsection{Identified Defects and Resolutions}
\begin{tabularx}{\linewidth}{lYY}\toprule Item & Finding & Resolution / evidence\\\midrule
BUILD-01 & JAR temporary replacement was denied by the sandbox. & Direct JAR temporary files into the writable build folder; subsequent build succeeded.\\
BUILD-02 & PowerShell split the unquoted Java system-property argument. & Quote the complete headless property; the test runner executed successfully.\\
UI-01 & Sign-in intro paragraph was clipped. & Correct preferred height and regenerate screenshots.\\
SEC-01 & Account APIs exposed credential records. & Return UserSummary; restrict directory access (TC-60--61).\\
INPUT-01 & Null credentials/role/winner and relative URLs caused poor failures. & Validate before mutation (TC-62--64, 70).\\
SEC-02 & Malformed password parameters and multiline CSV formulas. & Bound derivation parameters and escape formula cells (TC-65--66).\\
DATA-01 & Oversized saves could become unreadable; closed repositories could write. & Validate before replacement and enforce lock lifetime (TC-67--68).\\
BUILD-03 & Maven did not execute custom checks; shell line endings were fragile. & JUnit adapter, UI profile, LF shell and Git attributes; the POSIX script can run the same suites under a suitable shell.\\
UI-02 & Swing could interpret user names as HTML. & Plain-text labels, lists and tables; escaped wrapped text and tooltips (UI-16).\\
DATA-02 & Package and model changes could invalidate old saved files. & Preserve serialized record names; migrate completed and active legacy fixtures (TC-69, AR-25--26).\\
UI-03 & League team names had insufficient width and saving returned to Teams. & Allocate wider identity columns and retain the selected Fixtures tab (UI-21).\\
RULE-01 & Activity type and progression were mixed. & Independent strategies and service compatibility checks (AR-01--05).\\
DATA-03 & Saving a result can also create a new round. & Roll back both changes and the audit on failure (AR-22--24, 27).\\\bottomrule\end{tabularx}
No unresolved failure was observed in the executed service suite. This statement does not imply absence of defects outside its coverage. Full end-to-end desktop mouse/keyboard coverage, large-data responsiveness, power-loss recovery, and cross-platform execution remain validation gaps.

\clearpage\section{Conclusion}
The suite provides executable evidence for the central OOAD behaviors: format-specific policies, protected actions, lifecycle guards, and persistence integrity. The known limits are documented and the tests can be rerun locally without an external framework. Future testing should extend GUI interaction coverage, add randomized score sequences and large-data responsiveness testing before extending the application to real campus deployment.
\subsection*{References}
\begingroup\raggedright\begin{enumerate}
\item Object Management Group. \emph{Unified Modeling Language, version 2.5.1}. \url{https://www.omg.org/spec/UML/2.5.1}. Accessed 29 October 2026. Used as the notation reference for the UML models.
\item Oracle. \emph{Java SE 17: javax.swing package}. \url{https://docs.oracle.com/en/java/javase/17/docs/api/java.desktop/javax/swing/package-summary.html}. Accessed 29 October 2026. Used for desktop components and event-dispatch-thread conventions.
\item Oracle. \emph{Java SE 17: ObjectInputFilter}. \url{https://docs.oracle.com/en/java/javase/17/docs/api/java.base/java/io/ObjectInputFilter.html}. Accessed 29 October 2026. Used for bounded, filtered repository deserialization.
\item Project materials. Task outlines in \code{requirements/system-analysis-requirements.txt}, \code{requirements/system-design-requirements.txt}, \code{requirements/implementation-report-requirements.txt}, \code{requirements/testing-requirements.txt}, and \code{requirements/sample-data-requirements.txt}; instructor feedback supplied by the team. These define the coursework scope, not an independent evaluation of the implementation.
\end{enumerate}\endgroup

\end{document}
